% lp-insurance-demand.tex — Structural Econometric Specification
% LP Insurance Demand: Quadratic Fee Concentration Exit Model
% Derived via Reiss & Wolak (2007) framework
% Key result: inverted-U relationship between fee concentration and LP exits

\documentclass[11pt,a4paper]{article}

\usepackage{amsmath,amssymb,amsthm}
\usepackage{booktabs}
\usepackage{geometry}
\usepackage{hyperref}
\usepackage{cleveref}
\usepackage{enumitem}
\usepackage{longtable}
\usepackage{array}
\usepackage{graphicx}

\geometry{margin=1in}

\newtheorem{assumption}{Assumption}
\newtheorem{proposition}{Proposition}
\newtheorem{definition}{Definition}
\newtheorem{hypothesis}{Hypothesis}
\newtheorem{result}{Result}

\title{Structural Econometric Specification:\\[4pt]
       \large Fee Concentration Risk and LP Exit: An Inverted-U\\[2pt]
       \large Insurance Demand from the Quadratic Deviation Model}
\author{ThetaSwap}
\date{\today}

\begin{document}

\maketitle

\begin{abstract}
We estimate the effect of fee concentration risk on passive LP exit
probability in Uniswap~V3. Using per-position fee shares from on-chain
Collect events, we compute the real fee concentration index
$A_T = \sqrt{\sum_k \theta_k x_k^2}$ and test its deviation from the
Ma~\&~Crapis (2024) equal-share null $A_T^{1/N}$ as the treatment variable
in a position-day exit hazard model.

A linear specification yields $\beta_1 < 0$: higher concentration is
associated with \emph{fewer} exits (shelter effect). However, a quadratic
specification reveals a statistically significant inverted-U relationship:
$\beta_{\mathrm{lin}} = -23.18$ ($p = 0.012$),
$\beta_{\mathrm{quad}} = +129.20$ ($p = 0.030$), with a turning point at
deviation $\approx 0.09$. Below this threshold, fee concentration proxies
for pool health; above it, the Capponi~\&~Zhu (2024) exit mechanism
activates---passive LPs' effective fee rate drops below their exit
threshold.

This identifies a concrete insurance trigger for ThetaSwap: when
concentration exceeds 9\% above the competitive equilibrium, passive LP exit
risk increases and insurance demand is positive. Five papers from the AMM
literature independently validate each component of this finding.
\end{abstract}

\tableofcontents

% ======================================================================
\section{Research Question}\label{sec:research-question}
% ======================================================================

\subsection{Economic Question}

\textbf{Primary:} Does fee concentration risk ($A_T$) drive passive LP
exits in Uniswap~V3?

\textbf{Secondary:} At what concentration threshold does the exit effect
activate, and can it be priced as an insurance premium?

\textbf{Framing:} This is an insurance demand model. The goal is to identify
the regime where concentration risk becomes material and to translate the
effect into a premium for ThetaSwap's fee concentration derivative.

\subsection{Unit of Observation}

\textbf{Position-day panel.} One observation per position per day over a
41-day window (2025-12-05 to 2026-01-14). Binary outcome: $\text{exit}_{i,t}
\in \{0, 1\}$. Treatment: fee concentration deviation from the competitive
null, measured at the pool-day level with lag.

\subsection{Outcome Variable}

\textbf{Binary exit indicator:}
\begin{equation}\label{eq:outcome}
  Y_{i,t} = \begin{cases}
    1 & \text{if position } i \text{ burns on day } t \\
    0 & \text{otherwise}
  \end{cases}
\end{equation}

Panel size: 3{,}365 position-day observations, 597 exit events, 40 cluster
days, from 600 positions in the ETH/USDC 30bps pool.


% ======================================================================
\section{The Fee Concentration Index}\label{sec:a-t}
% ======================================================================

\begin{definition}[Fee Concentration Index]\label{def:a-t}
For a pool with $N$ active positions on day $t$, the fee concentration
index is:
\begin{equation}\label{eq:a-t}
  A_T = \sqrt{\sum_{k=1}^{N} \theta_k \, x_k^2}
\end{equation}
where $x_k = \text{fee}_k / \sum_j \text{fee}_j$ is position $k$'s fee
share and $\theta_k = 1/\text{blocklife}_k$ is its turnover rate.
\end{definition}

$A_T$ aggregates two dimensions of competitive pressure: \emph{who earns
the fees} ($x_k$) and \emph{how fast positions turn over} ($\theta_k$).
High $A_T$ means a few short-lived positions capture most fee revenue.

\subsection{The Equal-Share Null}

\begin{definition}[Ma--Crapis Null]\label{def:null}
Under the Ma~\&~Crapis (2024) symmetric equilibrium with $N$ identical LPs,
each earns equal fee share $x_k = 1/N$:
\begin{equation}\label{eq:null}
  A_T^{1/N} = \sqrt{\frac{\sum_k \theta_k}{N^2}}
\end{equation}
This is the competitive baseline---the value $A_T$ takes when fee
concentration is zero.
\end{definition}

\begin{definition}[Concentration Deviation]\label{def:deviation}
The treatment variable is the excess concentration above the null:
\begin{equation}\label{eq:deviation}
  \Delta_t \equiv A_{T,t}^{\text{real}} - A_{T,t}^{1/N}
\end{equation}
computed from the \textbf{same} set of positions on each day. $\Delta_t > 0$
indicates concentration above the competitive equilibrium; $\Delta_t < 0$
indicates more equal fee distribution than the symmetric benchmark.
\end{definition}

\subsection{Empirical Properties of $\Delta_t$}

Data source: Dune query 6783604, joining \texttt{Collect}, \texttt{Burn},
and \texttt{Mint} events for the ETH/USDC 30bps pool
(\texttt{0x8ad599c3...eb6e6d8}). Fee amounts isolated by subtracting Burn
principal from Collect amounts. Token ordering: token0 = USDC (6 dec),
token1 = WETH (18 dec).

\begin{center}
\begin{tabular}{ll}
\toprule
Statistic & Value \\
\midrule
Days with $\Delta > 0$ & 26/41 (63\%) \\
Mean $A_T^{\text{real}} / A_T^{1/N}$ & 2.65$\times$ \\
Median deviation & +0.0004 \\
Range & $[-0.032, +0.158]$ \\
Positions per day & 13--65 \\
\bottomrule
\end{tabular}
\end{center}

The real $A_T$ exceeds the null on most days, confirming empirically that
fee shares are asymmetric---consistent with Aquilina et al.\ (2024) finding
that 7\% of LP addresses capture 80\% of fees.


% ======================================================================
\section{Theoretical Foundations}\label{sec:theory}
% ======================================================================

Five papers independently validate the components of $A_T$ and the exit
mechanism.

\subsection{Capponi \& Zhu (2024): Optimal Exit Threshold}

\emph{``Optimal Exiting for Liquidity Provision in Constant Function Market
Makers.''}

The LP solves an optimal stopping problem with exit thresholds $p_U^*$ and
$p_L^*$. The key comparative static:
\begin{equation}\label{eq:capponi-zhu}
  \frac{\partial p_U^*}{\partial \phi} > 0
\end{equation}
A higher effective fee rate $\phi$ raises the exit threshold (LP stays
longer). When $A_T$ rises, passive LP's effective fee rate $\hat\phi_i =
x_i \cdot \phi_{\text{pool}}$ drops because $x_i < 1/N$. Lower $\hat\phi
\Rightarrow$ lower $p_U^* \Rightarrow$ earlier exit.

\textbf{Implication for $A_T$:} $A_T$ enters the exit decision through
$\hat\phi$. This predicts $\beta_1 > 0$ for sufficiently high concentration.

\subsection{Capponi, Jia \& Zhu (2024): JIT Fee Dilution}

\emph{``The Paradox of Just-in-Time Liquidity in Decentralized Exchanges.''}

JIT liquidity creates fee share asymmetry: the passive LP's share is
$d_P / (d_P + d_J)$ where $d_J$ is JIT depth. This is exactly the $x_k$
component of $A_T$. JIT entry \emph{mechanically} raises $A_T$ by
concentrating fees in short-lived positions.

\textbf{Implication for identification:} JIT activity is a natural
instrument for $A_T$. Our JIT control analysis (32 JIT positions, 5.3\%
share) shows JIT count has an independent positive effect on exits
($\beta_{\text{JIT}} = 0.177$, $p = 0.031$), but does not explain the
$A_T$ effect---confirming that $A_T$ captures concentration \emph{beyond}
JIT.

\subsection{Ma \& Crapis (2024): The Competitive Null}

\emph{``Decentralized Exchange Design: Permissionless Liquidity Provision.''}

$N$ symmetric LPs in free-entry equilibrium: each earns $\Pi_i =
(1/N^2) f B_D$. The equal-share outcome $x_k = 1/N$ is the \textbf{null
hypothesis} of $A_T$. Our deviation $\Delta_t = A_T^{\text{real}} -
A_T^{1/N}$ directly tests departure from this equilibrium.

\textbf{Implication for the model:} $\Delta_t = 0$ means the pool operates
at the competitive equilibrium. $\Delta_t > 0$ means fee concentration
exceeds what symmetric competition predicts. The deviation is the correct
treatment variable---not $A_T$ in levels.

\subsection{Aquilina, Foley, Gambacorta \& Krekel (2024): Empirical
Validation}

\emph{``Decentralised Dealers? Examining Liquidity Provision in
Decentralised Exchanges.'' BIS Working Paper 1227.}

Sophisticated LPs (7\% of addresses) provide 65--85\% of liquidity, use
tick ranges 3.4$\times$ tighter, hold positions 7.5$\times$ shorter, and
capture 75--90\% of fees. This validates that $x_k$ in $A_T$ is highly
skewed, not symmetric.

Their profitability decomposition $R_{\text{net}} = \text{FeeYield} +
\text{IL} - \text{Gas}$ shows that passive LPs earn negative returns when
sophisticated LPs dominate---precisely when $\Delta_t$ is large.

\subsection{Bichuch \& Feinstein (2024): Fee-as-Derivative Pricing}

\emph{``The Price of Liquidity: Implied Volatility of AMM Fees.''}

The LP value function is an optimal stopping problem with implied fee rate
$F(p_x, p_y) = p_y \cdot \ell(p_x/p_y)$ where $\ell$ is instantaneous LVR.
Theorem~4.2: a risk-neutral LP is indifferent across all stopping times.
This means \emph{any} deterministic exit driver must come from fee rate
heterogeneity across LPs---exactly what $A_T$ measures.

Their fixed-for-floating fee swap (Corollary~5.5) is structurally identical
to ThetaSwap's insurance product: the LP exchanges volatile realized fees
for a fixed premium.

\textbf{Implication:} $A_T$ measures the \emph{distribution} of the pool
fee rate across LPs. Bichuch \& Feinstein's framework prices the pool-level
fee; $A_T$ determines who receives it.


% ======================================================================
\section{Economic Model}\label{sec:economic-model}
% ======================================================================

\subsection{Environment and Actors}

Single Uniswap~V3 pool (ETH/USDC 30bps). Two agent types:

\begin{enumerate}
  \item \textbf{Passive LP (PLP):} Commits liquidity for extended periods.
        Observes pool state (reserves, price, accumulated fees). Does
        \emph{not} observe $A_T$ directly but reacts to its symptoms (lower
        fee revenue). Decision: when to exit.

  \item \textbf{JIT/Sophisticated LP:} Not explicitly modeled as strategic.
        Enters as a latent force that raises $A_T$ by capturing
        disproportionate fee shares. Characterized by short blocklife and
        tight tick ranges (Aquilina et al.).
\end{enumerate}

\subsection{Equilibrium Concept}

Competitive (price-taking) per Capponi \& Zhu (2024). The PLP takes fee
rate, price process, and concentration risk as given and solves an optimal
stopping problem.

\begin{assumption}[Price-Taking LP]\label{asm:price-taking}
The PLP solves:
\begin{equation}
  \max_{\tau} \; \mathbb{E}\left[\int_0^\tau e^{-rt}
  \hat\phi_i \, dF_t - dL_t \;\middle|\; \mathcal{F}_0\right]
\end{equation}
where $\hat\phi_i = x_i \cdot \phi_{\text{pool}}$ is the LP's effective fee
rate given fee share $x_i$, $F_t$ is cumulative fee revenue, and $L_t$ is
impermanent loss.
\end{assumption}


% ======================================================================
\section{Stochastic Model}\label{sec:stochastic-model}
% ======================================================================

\subsection{Unobserved Heterogeneity ($\eta$)}

\begin{enumerate}
  \item \textbf{LP risk aversion} ($\eta_i^{\text{risk}}$): heterogeneous
        preferences. Some PLPs tolerate more concentration. NFT Position
        Manager masks identity.

  \item \textbf{Outside options} ($\eta_i^{\text{outside}}$): competing pool
        returns, staking yields, CEX opportunities.
\end{enumerate}

Bias direction: \textbf{toward zero} (conservative). Risk-averse LPs avoid
entering during high-$A_T$ periods \emph{and} exit faster, attenuating
$\beta$.

\subsection{Agent Uncertainty}

Future price process only ($u_t \sim \text{GBM}$, per Capponi \& Zhu).
Drives IL realization.

\subsection{Implied Error Structure}

\begin{equation}\label{eq:error-structure}
  \varepsilon_i = \underbrace{\eta_i^{\text{risk}} +
  \eta_i^{\text{outside}}}_{\text{unobserved heterogeneity}}
  + \underbrace{\nu_i}_{\text{measurement error}}
\end{equation}


% ======================================================================
\section{Estimation Strategy}\label{sec:estimation}
% ======================================================================

\subsection{Model Progression}

We estimate three nested specifications to motivate the quadratic model.

\subsubsection{Specification 1: Linear (Point Lag)}

\begin{equation}\label{eq:linear}
  P(\text{exit}_{i,t} = 1) = \sigma\!\left(
    \beta_0 + \beta_1 A_{T,t-\ell} + \beta_2 \, \text{IL}_t
    + \beta_3 \log(\text{age}_{i,t})
  \right)
\end{equation}

\begin{result}[Linear A\_T]\label{res:linear}
$\beta_1 = -3.93$, cluster $p = 0.072$. Marginally significant but
\textbf{wrong sign}---higher concentration associated with fewer exits.
Robust across lags 1--7 (all negative, 4/5 significant).
\end{result}

\subsubsection{Specification 2: Linear Deviation from 1/N}

\begin{equation}\label{eq:deviation-linear}
  P(\text{exit}_{i,t} = 1) = \sigma\!\left(
    \beta_0 + \beta_1 \Delta_{t-\ell} + \beta_2 \, \text{IL}_t
    + \beta_3 \log(\text{age}_{i,t})
  \right)
\end{equation}

where $\Delta_{t} = A_T^{\text{real}} - A_T^{1/N}$ from the same position
set.

\begin{result}[Linear Deviation]\label{res:deviation}
$\beta_1 = -5.44$, cluster $p = 0.109$. Still negative. But
\textbf{dose-response analysis} reveals non-monotonicity:

\begin{center}
\begin{tabular}{cccc}
\toprule
Quartile & Mean $\Delta$ & Exit Rate & Pattern \\
\midrule
Q1 (below null) & $-0.009$ & 17.75\% & baseline \\
Q2 (at null) & $-0.000$ & 17.87\% & baseline \\
\textbf{Q3 (mild conc.)} & $\mathbf{+0.002}$ & $\mathbf{20.56\%}$ &
  \textbf{exits increase} \\
Q4 (extreme conc.) & $+0.025$ & 14.37\% & exits decrease \\
\bottomrule
\end{tabular}
\end{center}

Q3 shows the predicted Capponi \& Zhu pattern. Q4 reverses---motivating
the quadratic specification.
\end{result}

\subsubsection{Specification 3: Quadratic Deviation (Primary Model)}

\begin{equation}\label{eq:quadratic}
\boxed{
  P(\text{exit}_{i,t} = 1) = \sigma\!\left(
    \beta_0 + \beta_1 \Delta_{t-\ell} + \beta_2 \Delta_{t-\ell}^2
    + \beta_3 \, \text{IL}_t + \beta_4 \log(\text{age}_{i,t})
  \right)
}
\end{equation}

This is the \textbf{primary specification}. The quadratic term captures the
inverted-U: shelter at low concentration, Capponi exit at high
concentration.


\subsection{Estimation}

Logit MLE via Newton--Raphson (IRLS) with \textbf{day-clustered sandwich
standard errors} (41 clusters):

\begin{equation}\label{eq:cluster-se}
  \widehat{\text{Var}}(\hat\beta) = H^{-1}
  \left(\sum_{g=1}^{G} S_g S_g'\right)
  H^{-1} \cdot \frac{G}{G-1} \cdot \frac{n}{n-k}
\end{equation}

where $S_g = \sum_{i \in g} (y_i - \hat p_i) x_i$ is the cluster-level
score and $H$ is the Hessian.

Implemented in JAX~0.9.1 CPU. Code: \texttt{econometrics/hazard.py}.


% ======================================================================
\section{Results}\label{sec:results}
% ======================================================================

\subsection{Main Result: Quadratic Deviation Model}

\begin{result}[Inverted-U]\label{res:quadratic}
Both terms of the quadratic deviation model are statistically significant at
the economically relevant lag horizon (1--3 days):

\begin{center}
\begin{tabular}{cccccc}
\toprule
Lag & $\hat\beta_1$ (linear) & $p$ & $\hat\beta_2$ (quadratic) & $p$ &
  Turning point \\
\midrule
\textbf{1} & $\mathbf{-23.18}$ & $\mathbf{0.012}$ & $\mathbf{+129.20}$ &
  $\mathbf{0.030}$ & $\mathbf{0.090}$ \\
2 & $-43.42$ & 0.016 & $+226.92$ & 0.065 & 0.096 \\
\textbf{3} & $\mathbf{-32.44}$ & $\mathbf{0.001}$ & $\mathbf{+205.34}$ &
  $\mathbf{0.004}$ & $\mathbf{0.079}$ \\
5 & $+3.81$ & 0.830 & $-76.31$ & 0.470 & --- \\
7 & $-5.55$ & 0.730 & $-16.46$ & 0.900 & --- \\
\bottomrule
\end{tabular}
\end{center}

The turning point clusters at $\Delta^* \approx 0.08\text{--}0.10$. The
signal dissipates at lags 5--7, consistent with PLPs responding to recent
concentration shocks, not stale ones.
\end{result}

\subsection{Full Lag-1 Coefficients}

\begin{center}
\begin{tabular}{lrrl}
\toprule
Parameter & Estimate & Cluster SE & Interpretation \\
\midrule
$\beta_0$ (intercept) & varies & --- & baseline exit probability \\
$\beta_1$ ($\Delta$) & $-23.18$ & 9.21 & shelter below turning point \\
$\beta_2$ ($\Delta^2$) & $+129.20$ & 59.69 & exit above turning point \\
$\beta_3$ (IL) & $+13.52$ & --- & higher IL $\to$ more exits \\
$\beta_4$ ($\log$ age) & $-0.42$ & --- & older positions less likely to exit \\
\bottomrule
\end{tabular}
\end{center}

$n = 3{,}365$; exits = 597; clusters = 40; pseudo-$R^2 = 0.350$; mean
$P(\text{exit}) = 0.177$.


% ======================================================================
\section{Economic Interpretation}\label{sec:interpretation}
% ======================================================================

\subsection{The Inverted-U: Two Regimes}

The quadratic model identifies two regimes separated by the turning point
$\Delta^* \approx 0.09$:

\begin{enumerate}
  \item \textbf{Shelter regime} ($\Delta < \Delta^*$): Mild fee
        concentration correlates with high trading volume. All LPs
        benefit from elevated fee revenue, even if shares are unequal.
        The marginal effect $\partial P / \partial \Delta < 0$: more
        concentration, fewer exits.

  \item \textbf{Capponi regime} ($\Delta > \Delta^*$): Extreme
        concentration means passive LPs' effective fee rate $\hat\phi_i$
        drops below their exit threshold. Per Capponi \& Zhu
        \eqref{eq:capponi-zhu}: $\partial p_U^* / \partial \phi > 0$, so
        lower $\phi \Rightarrow$ lower $p_U^* \Rightarrow$ earlier exit.
        The marginal effect flips: $\partial P / \partial \Delta > 0$.
\end{enumerate}

\subsection{Marginal Effect at the Turning Point}

The marginal effect of the quadratic model is:
\begin{equation}\label{eq:marginal}
  \frac{\partial P(\text{exit})}{\partial \Delta} =
  \left(\beta_1 + 2\beta_2 \Delta\right) \cdot \bar P (1 - \bar P)
\end{equation}

At the turning point $\Delta^* = -\beta_1 / (2\beta_2) = 0.090$:
\[
  \frac{\partial P}{\partial \Delta}\bigg|_{\Delta^*} = 0
\]

For $\Delta = 0.15$ (well into the Capponi regime):
\begin{align}
  \beta_1 + 2\beta_2(0.15) &= -23.18 + 2(129.20)(0.15) = 15.58 \\
  \frac{\partial P}{\partial \Delta} &= 15.58 \times 0.177 \times 0.823
  = 2.27
\end{align}

A 0.01 increase in $\Delta$ at this level raises exit probability by
$\sim$2.3 percentage points.

\subsection{Insurance Pricing Translation}

For a passive LP experiencing $\Delta = 0.15$ (concentration 15\% above the
competitive null):

\begin{center}
\begin{tabular}{lr}
\toprule
Quantity & Value \\
\midrule
$\Delta P(\text{exit})$ per 0.01 $\Delta$ & +2.3 pp \\
Average remaining position hours & 48 \\
Hours of fee revenue at risk & $0.023 \times 48 = 1.1$ hrs \\
Fee revenue per hour (pool avg) & \$100 \\
Implied premium per position & \$110 \\
\bottomrule
\end{tabular}
\end{center}

This is the maximum willingness-to-pay for ThetaSwap's insurance: a PLP
would pay up to \$110 to eliminate the excess exit risk from concentration
above the turning point.


% ======================================================================
\section{Specification Tests}\label{sec:spec-tests}
% ======================================================================

\subsection{Test 1: Sign Restriction (Capponi \& Zhu Prediction)}

\begin{hypothesis}
$\beta_2 > 0$: the quadratic term is positive, indicating that extreme
concentration drives exits.
\end{hypothesis}

\textbf{Result: PASS.} $\hat\beta_2 = +129.20$, $p = 0.030$ at lag 1;
$+205.34$, $p = 0.004$ at lag 3. The Capponi exit mechanism is confirmed
in the high-concentration regime.

\subsection{Test 2: Dose-Response Monotonicity}

\begin{hypothesis}
Exit rates increase with $\Delta$ in the upper quartiles.
\end{hypothesis}

\textbf{Result: PASS (inverted-U).} Q1--Q2 stable (17.8\%), Q3 peak
(20.6\%), Q4 drops (14.4\%). The non-monotonicity motivates the quadratic
and is captured by the model.

\subsection{Test 3: Lag Robustness}

\begin{hypothesis}
Both $\beta_1$ and $\beta_2$ are significant for at least 3 of 5 lag
windows.
\end{hypothesis}

\textbf{Result: PASS.} Lags 1, 2, 3 all have significant quadratic terms
($p < 0.07$). The turning point is stable at 0.08--0.10. Signal dissipates
at lags 5--7, consistent with short-memory PLP reaction.

\subsection{Test 4: IL Orthogonality}

\begin{hypothesis}
$\beta_1, \beta_2$ persist with and without the IL control.
\end{hypothesis}

\textbf{Result: PASS.} IL coefficient ($\beta_3 = 13.52$) is positive
(higher IL $\to$ more exits) but does not absorb the concentration effect.
The quadratic structure is present in the $A_T$-only model as well.

\subsection{Test 5: JIT Confounding Control}

\begin{hypothesis}
The concentration effect is not an artifact of JIT activity.
\end{hypothesis}

\textbf{Result: PASS.} Adding a JIT proxy (positions with blocklife $< 100$
blocks) as a control:
\begin{itemize}[nosep]
  \item JIT count has independent positive effect on exits
        ($\beta_{\text{JIT}} = 0.177$, $p = 0.031$)
  \item The $A_T$ effect becomes \emph{more} significant ($p: 0.072 \to
        0.022$) with JIT controlled
  \item Confirms: $A_T$ captures concentration \emph{beyond} JIT
\end{itemize}

\subsection{Test 6: Alternative Treatment Timing}

Three treatment constructions tested:

\begin{center}
\begin{tabular}{lccl}
\toprule
Treatment & $\beta_1$ (linear) & $p$ & Finding \\
\midrule
Point A\_T (real, lag 1) & $-3.93$ & 0.072 & Negative (shelter) \\
Lifetime mean A\_T & $-7.60$ & 0.444 & Negative, insignificant \\
Deviation from 1/N & $-5.44$ & 0.109 & Negative (Q3 reveals non-linearity) \\
\textbf{Quadratic deviation} & \multicolumn{2}{c}{\textbf{See Result~3}} &
  \textbf{Inverted-U confirmed} \\
\bottomrule
\end{tabular}
\end{center}


% ======================================================================
\section{Robustness: Five Linear Specifications}\label{sec:robustness}
% ======================================================================

Before arriving at the quadratic model, we tested five linear specifications
that progressively refined the treatment variable. All failed to produce a
positive $\beta_1$, but each contributed to understanding the data-generating
process.

\subsection{Specification A: Point Real A\_T}

Treatment = $A_{T,t-\ell}^{\text{real}}$ (single day's concentration).

\begin{center}
\begin{tabular}{ccccc}
\toprule
Lag & $\hat\beta_1$ & Cluster SE & $p$ & Sig \\
\midrule
1 & $-3.928$ & 2.186 & 0.072 & * \\
2 & $-8.783$ & 2.227 & 0.0001 & *** \\
3 & $-0.414$ & 1.131 & 0.714 & \\
5 & $-6.523$ & 1.650 & 0.0001 & *** \\
7 & $-5.625$ & 2.261 & 0.013 & ** \\
\bottomrule
\end{tabular}
\end{center}

\textbf{Finding:} Consistently negative. A naive interpretation is that
concentration \emph{protects} LPs. But the proxy A\_T (x\_k = 1/N) showed the
same pattern ($\beta_1 = -4.90$, $p = 0.15$), suggesting the level of
$A_T$ confounds pool health (volume) with concentration pressure.

\subsection{Specification B: Lifetime Mean A\_T}

Treatment = mean $A_T^{\text{real}}$ over $[\text{mint}_i, t - \ell]$.
Motivated by the hypothesis that passive LPs react to accumulated exposure,
not point shocks.

\begin{center}
\begin{tabular}{ccccc}
\toprule
Lag & $\hat\beta_1$ & Cluster SE & $p$ & Sig \\
\midrule
1 & $-7.595$ & 9.913 & 0.444 & \\
3 & $-11.603$ & 7.959 & 0.145 & \\
5 & $-14.977$ & 8.021 & 0.062 & * \\
7 & $-3.042$ & 6.829 & 0.656 & \\
\bottomrule
\end{tabular}
\end{center}

\textbf{Finding:} Amplifies the negative effect but with much larger SEs
(reduced variation in running means). Marginal significance only at lag 5.

\subsection{Specification C: Deviation from 1/N (Wrong Null)}

First attempt at the deviation treatment, using the old proxy A\_T from a
\emph{different} position set (Q4 burns) as the null. The real and null
A\_T were computed from different populations---an apples-to-oranges
comparison.

\textbf{Finding:} Real A\_T was 10--100$\times$ smaller than the proxy null
(different position sets, different blocklives). Deviation was negative for
40/41 days. Only lag-1 showed positive $\beta_1$ (+1.92, $p = 0.43$).
\textbf{Discarded} after identifying the data mismatch.

\subsection{Specification D: Corrected Deviation (Same Position Set)}

Null A\_T recomputed from the \emph{same} Dune query (same positions, same
blocklives). Now 26/41 days have $\Delta > 0$, mean ratio = 2.65$\times$.

\textbf{Finding:} Still negative $\beta_1$ in linear model, but
dose-response revealed the Q3 peak (20.56\%) and Q4 dip (14.37\%) that
motivated the quadratic specification.

\subsection{Specification E: JIT-Controlled Linear}

Added daily JIT count (positions with blocklife $< 100$) as control.

\textbf{Finding:} $\beta_1$ became \emph{more} negative and more significant
($p: 0.072 \to 0.022$). JIT itself predicts exits but does not confound the
$A_T$ effect.


% ======================================================================
\section{Literature Alignment}\label{sec:literature}
% ======================================================================

Each component of the quadratic deviation model is validated by at least one
paper:

\begin{longtable}{p{3.5cm} p{4cm} p{5cm}}
\toprule
Paper & Component Validated & Connection to Our Result \\
\midrule
\endhead

Capponi \& Zhu (2024) &
  Exit threshold: $\partial p_U^* / \partial \phi > 0$ &
  Confirmed in the Capponi regime ($\Delta > 0.09$): $\beta_2 > 0$ means
  extreme concentration drives exits \\[6pt]

Capponi, Jia \& Zhu (2024) &
  $x_k$ captures JIT fee dilution &
  JIT control shows $A_T$ effect is independent of JIT confounding; $x_k$
  reflects broader concentration than just JIT \\[6pt]

Ma \& Crapis (2024) &
  $x_k = 1/N$ is the competitive null &
  Deviation $\Delta = A_T^{\text{real}} - A_T^{1/N}$ is the correct
  treatment; levels confound pool health \\[6pt]

Aquilina et al.\ (2024) &
  7\% of LPs capture 80\% of fees &
  Real $A_T$ is 2.65$\times$ the null on average, confirming massive fee
  share asymmetry \\[6pt]

Bichuch \& Feinstein (2024) &
  Risk-neutral LP indifferent on stopping times &
  Any exit driver must come from fee \emph{heterogeneity} across LPs, not
  pool-level fee rate---exactly what $\Delta$ measures \\

\bottomrule
\end{longtable}


% ======================================================================
\section{Connection to ThetaSwap}\label{sec:thetaswap}
% ======================================================================

\begin{enumerate}
  \item \textbf{Insurance trigger:} The turning point $\Delta^* \approx
        0.09$ defines when insurance becomes valuable. When concentration
        exceeds 9\% above the Ma--Crapis equilibrium, passive LP exit risk
        increases. ThetaSwap monitors $\Delta_t$ in real time and activates
        payouts when $\Delta > \Delta^*$.

  \item \textbf{Premium calibration:} The quadratic marginal effect at
        any $\Delta$ translates directly to an insurance premium via
        \cref{eq:marginal}. At $\Delta = 0.15$: \$110 per position per
        event.

  \item \textbf{Distinct from IL hedging:} The IL coefficient ($\beta_3$)
        is positive but does not absorb the concentration effect. $A_T$
        captures a risk orthogonal to price-movement risk, confirming that
        ThetaSwap addresses a gap unmet by options or IL protection products.

  \item \textbf{Product structure:} Per Bichuch \& Feinstein's
        fixed-for-floating fee swap framework, ThetaSwap is a derivative
        that pays the PLP when realized fee share $x_i$ falls below $1/N$
        by more than $\Delta^*$. The premium is the implied volatility of
        fee shares.

  \item \textbf{Market sizing:} 63\% of days in our sample have
        $\Delta > 0$. Of these, days exceeding $\Delta^*$ represent the
        insurance-triggering events. Aggregating the implied premium across
        600 positions gives the minimum viable insurance pool size.
\end{enumerate}


% ======================================================================
\section{Data and Implementation}\label{sec:data}
% ======================================================================

\begin{table}[h]
\centering
\begin{tabular}{llll}
\toprule
Query & Description & Rows & Cost \\
\midrule
Q4v2 & Position-level burns, blocklife & 600 & 0.18 cr \\
Q5 & Daily IL proxy & 41 & 0.14 cr \\
Q6 & Per-position Collect fees + blocklife & 41 (aggregated) & 3.3 cr \\
 & (real $A_T$ and null $A_T$ per day) & & \\
\bottomrule
\end{tabular}
\caption{Dune MCP queries. Total cost: $\sim$3.6 credits.}
\end{table}

\textbf{Pool:} ETH/USDC 30bps
(\texttt{0x8ad599c3a0ff1de082011efddc58f1908eb6e6d8}).

\textbf{Window:} 2025-12-05 to 2026-01-14 (41 days).

\textbf{Implementation:} Python 3.11, JAX 0.9.1 CPU,
\texttt{@functional-python} patterns (frozen dataclasses, pure functions).
Code in \texttt{econometrics/\{types,data,ingest,hazard\}.py}. 52 unit
tests passing.

Key functions:
\begin{itemize}[nosep]
  \item \texttt{build\_exit\_panel\_deviation()}: constructs position-day
        panel with $\Delta$ as treatment
  \item \texttt{logit\_mle\_quadratic()}: Newton--Raphson logit with
        $\Delta + \Delta^2$ and clustered SEs
  \item \texttt{marginal\_effect\_at\_means()}: translates $\beta$ to
        insurance premium
\end{itemize}


% ======================================================================
\section{References}\label{sec:references}
% ======================================================================

\begin{itemize}[nosep]
  \item Aquilina, M., Foley, S., Gambacorta, L., \& Krekel, M.\ (2024).
        ``Decentralised Dealers? Examining Liquidity Provision in
        Decentralised Exchanges.'' \textit{BIS Working Paper} No.~1227.

  \item Bichuch, M.\ \& Feinstein, Z.\ (2024). ``The Price of Liquidity:
        Implied Volatility of AMM Fees.''

  \item Capponi, A.\ \& Zhu, B.\ (2024). ``Optimal Exiting for Liquidity
        Provision in Constant Function Market Makers.''

  \item Capponi, A., Jia, R., \& Zhu, B.\ (2024). ``The Paradox of
        Just-in-Time Liquidity in Decentralized Exchanges: More Providers
        Can Lead to Less Liquidity.''

  \item Ma, J.\ \& Crapis, H.\ (2024). ``Decentralized Exchange Design:
        Permissionless Liquidity Provision.''

  \item Reiss, P.C.\ \& Wolak, F.A.\ (2007). ``Structural Econometric
        Modeling: Rationales and Examples from Industrial Organization.''
        \textit{Handbook of Econometrics}, Vol.~6A, Ch.~64.
\end{itemize}

\end{document}
